% ECONOMICS_PROBLEM: {"id": "I26-514", "title": "Returns to education: mechanisms and heterogeneity", "jel_code": "I26", "source_id": "OP-0021"}
% !TeX program = xelatex
\documentclass[11pt,a4paper]{article}
\usepackage[margin=23mm]{geometry}
\usepackage{fontspec}
\setmainfont{Latin Modern Roman}
\usepackage{amsmath,amssymb,mathtools}
\usepackage{xcolor,enumitem,booktabs,longtable,array,xurl,hyperref}
\definecolor{muted}{HTML}{53636C}
\hypersetup{colorlinks=true,urlcolor=blue,linkcolor=blue}
\setlength{\parindent}{0pt}
\setlength{\parskip}{5pt}
\newcommand{\R}{\mathbb R}
\newcommand{\E}{\mathbb E}
\newcommand{\N}{\mathbb N}
\newcommand{\ind}{\mathbf 1}
\newcommand{\Var}{\operatorname{Var}}
\newcommand{\Cov}{\operatorname{Cov}}
\newcommand{\supp}{\operatorname{supp}}
\DeclareMathOperator*{\argmin}{arg\,min}
\DeclareMathOperator*{\argmax}{arg\,max}
\newcommand{\entrymeta}[3]{{\sffamily\small\color{muted}\textbf{\texttt{#1}}\quad|\quad #2\par #3\par}}
\newcommand{\Problemref}[1]{\texttt{#1}}
\newcommand{\JELref}[2]{\texttt{#2} (JEL \texttt{#1})}
\begin{document}
\section*{Returns to education: mechanisms and heterogeneity}
\textbf{Problem ID:} \texttt{I26-514}\quad \textbf{JEL:} \texttt{I26}\par
% BEGIN ECONOMICS STATEMENT
\label{problem:OP-0021}
\label{atlas:prob:L1}
\noindent\textbf{Original atlas ID:} L1 (entry 21). \textbf{Classification:} Editorial causal-mechanism identification and transport problem.

\noindent\textbf{Original atlas question:} How much of education’s return reflects skills, signaling, networks or selection, and for whom?

\subsection*{Definitions and assumptions}
Fix a population, a schooling margin, and an observation horizon. Let $D\in\{0,1\}$ denote completing the specified educational program, $X$ denote pretreatment characteristics, and $Y_i(d)$ denote individual $i$'s discounted lifetime earnings under schooling status $d$, with prices and other people's schooling held fixed. Assume integrability and consistency. Program identity includes institution, field, duration, and credential; one additional year is not automatically the same treatment in different programs. Define the conditional earnings effect $\tau(x)=\mathbb E[Y_i(1)-Y_i(0)\mid X_i=x]$.

For a mechanism decomposition, specify a modular structural causal model with channels $K=\{h,c,n,r\}$: skills, credential information, networks, and remaining direct pathways. Let $Y_i(S)$ be earnings when exactly the channels in $S\subseteq K$ transmit the schooling change, according to explicitly defined channel interventions. Require $Y_i(\varnothing)=Y_i(0)$ and $Y_i(K)=Y_i(1)$. Such interventions are additional model assumptions, not quantities defined by observing a mediator.
\subsection*{Formal research question}
For each channel $k$, identify or sharply bound
\[
 \phi_k(x)=\sum_{S\subseteq K\setminus\{k\}}
 \frac{|S|!(|K|-|S|-1)!}{|K|!}
 \mathbb E[Y_i(S\cup\{k\})-Y_i(S)\mid X_i=x].
\]
These Shapley allocations satisfy $\sum_{k\in K}\phi_k(x)=\tau(x)$ and allocate interaction effects by an explicit convention. Given an observed law $P$ and an announced model class $\mathcal M$, the identification target is the set of $(\tau,\phi)$ generated by models in $\mathcal M$ consistent with $P$. Determine which additional experiments or restrictions reduce this set and permit transport to a second population.
\subsection*{Interpretation and scope}
Selection is a difference in who takes schooling, not a causal transmission channel. An instrumental variable identifies a complier effect only under independence, relevance, exclusion, and monotonicity, and generally does not identify the displayed mechanism allocation. Mediator measurements need intervention or structural restrictions. A causal earnings effect is also different from an internal rate of return: tuition, foregone earnings, taxes, uncertainty, and discounting must enter an investment calculation. General-equilibrium credential dilution requires a separate enrollment-policy model.
\subsection*{Source audit and status}
Sources [S1]--[S3] verify the cited education literature and already provide substantial causal and conceptual results. The joint mechanism-and-transport target above is an editorial formulation. Neither the bibliography nor a local schooling estimate establishes a universally unresolved decomposition, or certifies its status throughout 2026.

\subsection*{References}
\begin{enumerate}[label=\textnormal{[S\arabic*]},leftmargin=*]
\item Joshua D. Angrist and Alan B. Krueger (1991). \emph{Does Compulsory School Attendance Affect Schooling and Earnings?}. The Quarterly Journal of Economics 106(4), 979--1014. \url{https://doi.org/10.2307/2937954}. \textit{Locator:} Primary publisher, working-paper, or author record: bibliographic information and abstract; full-text theorem verification is not claimed. \textit{Role:} Compulsory-schooling instrumental-variable evidence; motivates the distinction between local returns and mechanisms.
\item David Card (1999). \emph{The Causal Effect of Education on Earnings}. Handbook of Labor Economics 3A, chapter 30, 1801--1863. \url{https://doi.org/10.1016/S1573-4463(99)03011-4}. \textit{Locator:} Primary publisher, working-paper, or author record: bibliographic information and abstract; full-text theorem verification is not claimed. \textit{Role:} Survey of causal schooling effects and heterogeneous returns; background for extrapolation.
\item James J. Heckman, Lance J. Lochner, and Petra E. Todd (2006). \emph{Earnings Functions, Rates of Return and Treatment Effects: The Mincer Equation and Beyond}. Handbook of the Economics of Education 1, chapter 7, 307--458. \url{https://doi.org/10.1016/S1574-0692(06)01007-5}. \textit{Locator:} Primary publisher, working-paper, or author record: bibliographic information and abstract; full-text theorem verification is not claimed. \textit{Role:} Distinguishes earnings effects from investment rates of return; conceptual source for the mathematical correction.
\end{enumerate}

\subsection*{Common conventions: Source mathematical conventions}

Unless an entry states otherwise, agent and firm sets are finite. State and action spaces are measurable, strategies and outcome maps are measurable, probability laws are specified, and expectations exist for the targets under discussion. Infinite-horizon discount factors lie in $(0,1)$ and discounted objectives are finite. Norms, tolerances and complexity input sizes must be specified for computational comparisons. Constants or primitives that appear locally are inputs, not empirical facts asserted by this catalogue.

\textbf{Identification.} A statistical model $m\in\mathcal M$ implies an observable law $P_m$ and target $\theta(m)$. At law $P$, its identified set is
\[
\Theta(P;\mathcal M)=\{\theta(m):m\in\mathcal M,\ P_m=P\}.
\]
Point identification holds when this set is a singleton, or equivalently when $P_{m_1}=P_{m_2}$ implies $\theta(m_1)=\theta(m_2)$ for all compatible models. Sharp bounds mean bounds attained, or approached when the boundary is not attained, by compatible models. Writing this set is a definition, not a solution: the substantive task is to establish informative restrictions and derive or estimate the set. An unrestricted-confounding model may give only trivial bounds.

\textbf{Inference.} Identification, estimability and valid uncertainty quantification are separate questions. Fix $\alpha\in(0,1)$. A confidence procedure $C_n$ over a declared law class $\mathcal P$ has asymptotic uniform coverage at level $1-\alpha$ when
\[
\liminf_{n\to\infty}\inf_{P\in\mathcal P}\Pr_P\{\theta(P)\in C_n\}\geq1-\alpha.
\]
For set-identified targets the coverage event must be defined separately, for example coverage of the full identified set. If an entry uses identification-indexed models instead of uniquely indexed laws, the same coverage convention is applied over those models. Pointwise limits do not establish uniform validity, and asymptotic validity does not guarantee finite-sample precision. Sample sizes, dependence structures and weak-identification sequences are part of each application.

\textbf{Counterfactuals.} $Y_i(a)$ is a potential outcome under a specified intervention, including its timing, implementation and versions. It is not $Y_i$ conditional on voluntarily choosing $a$. A full intervention vector permits interference; shorthand scalar treatment notation does not silently impose no interference. Assignment, selection, attrition, anticipation and measurement must be specified. A claim about an instrument's local effect is not automatically a population policy effect.

\textbf{Economic equilibrium.} A counterfactual model includes best responses, constraints, feasibility, market clearing, state transitions and expectations consistent with the stated information. When several equilibria exist, either a declared selector is an input or the target is set-valued. Comparisons across policy regimes must use the stated selection convention. Reduced-form relationships are not presumed invariant after an equilibrium-changing intervention.

\textbf{Optimization and welfare.} Optimization is over the stated feasible set. If attainment is not established, $\sup$ or $\inf$ and $\varepsilon$-optimal choices replace an asserted maximizer or minimizer. For example, with value $V=\sup_{a\in\mathcal A}W(a)<\infty$, an $\varepsilon$-optimal set is $\{a\in\mathcal A:W(a)\geq V-\varepsilon\}$ for $\varepsilon>0$. Benefits, costs, risk and health or environmental outcomes can be aggregated only using declared units and conversion weights. Interpersonal utility weights, the treatment of fiscal transfers and foreign profits, discounting, and the behavioral welfare criterion are explicit inputs. A comparative-static derivative additionally requires existence and appropriate differentiability of the selected counterfactual.

\textbf{Sources.} Local labels [S1], [S2], and so forth restart in every question. Locators identify the relevant abstract, model, section or result at the level actually checked; a bibliographic or abstract-level verification is not represented as inspection of an entire proof. Working papers are distinguished from later publications and changing versions. Source summaries are paraphrased. All URL checks and editorial formulations refer to the audit date stated above.
% END ECONOMICS STATEMENT
\end{document}
